\documentclass[12pt,letterpaper]{article}

% =========================================================
% IDIOMA Y CODIFICACIÓN
% =========================================================

\usepackage[utf8]{inputenc}
\usepackage[T1]{fontenc}
\usepackage[spanish,es-nodecimaldot]{babel}

% =========================================================
% TIPOGRAFÍA
% Similar a Times New Roman
% =========================================================

\usepackage{newtxtext}
\usepackage{newtxmath}

% =========================================================
% MÁRGENES
% =========================================================

\usepackage[
    letterpaper,
    top=2.5cm,
    bottom=2.5cm,
    left=3cm,
    right=3cm,
    footskip=1.25cm
]{geometry}

% Si quieres VER los márgenes mientras editas,
% descomenta temporalmente la siguiente línea:
%
% \usepackage{showframe}
%
% NO debe quedar activa en la versión final.

% =========================================================
% INTERLINEADO Y PÁRRAFOS
% =========================================================

\usepackage{setspace}
\onehalfspacing

\setlength{\parindent}{1.25cm}
\setlength{\parskip}{0pt}

% =========================================================
% JUSTIFICACIÓN
% =========================================================

\usepackage{ragged2e}

% =========================================================
% PIE DE PÁGINA
% =========================================================

\usepackage{fancyhdr}

\pagestyle{fancy}
\fancyhf{}

\fancyfoot[L]{\small FES Acatlán · UNAM}
\fancyfoot[R]{\small \thepage}

\renewcommand{\headrulewidth}{0pt}
\renewcommand{\footrulewidth}{0.4pt}

% =========================================================
% FORMATO DE SECCIONES
% =========================================================

\usepackage{titlesec}

\titleformat{\section}
    {\normalfont\bfseries\large}
    {\thesection.}
    {0.6em}
    {}

\titleformat{\subsection}
    {\normalfont\bfseries\normalsize}
    {\thesubsection.}
    {0.6em}
    {}

\titlespacing*{\section}
    {0pt}
    {18pt}
    {12pt}

\titlespacing*{\subsection}
    {0pt}
    {15pt}
    {9pt}

% =========================================================
% ÍNDICE
% =========================================================

\usepackage{tocloft}

\renewcommand{\contentsname}{Índice}

\renewcommand{\cftsecfont}{\normalfont}
\renewcommand{\cftsecpagefont}{\normalfont}

% =========================================================
% HIPERVÍNCULOS
% =========================================================


\usepackage[hidelinks]{hyperref}

\hypersetup{
    pdftitle={
        Contenido sintético en Internet y degradación de modelos de lenguaje
    },
    pdfsubject={
        Investigación en Matemáticas Aplicadas y Computación
    },
    pdfauthor={
        Carlos Emiliano Segura Loera,
        Valeria Hernandez Santos,
        Nehiel Ramirez Hernandez
    },
    pdfkeywords={
        inteligencia artificial generativa,
        modelos de lenguaje,
        contenido sintético,
        model collapse,
        datos sintéticos,
        diversidad lingüística
    }
}

% =========================================================
% MEJORAS TIPOGRÁFICAS
% =========================================================

\usepackage{microtype}
\raggedbottom

% =========================================================
% DOCUMENTO
% =========================================================

\begin{document}

% =========================================================
% PORTADA
% =========================================================

\begin{titlepage}

\thispagestyle{empty}

% Interlineado sencillo exclusivamente en la portada
\begin{spacing}{1}

% ---------------------------------------------------------
% INSTITUCIÓN
% ---------------------------------------------------------

\begin{center}

\vspace*{0.5cm}

{\large\bfseries
UNIVERSIDAD NACIONAL AUTÓNOMA DE MÉXICO
}

\vspace{0.25cm}

{\large
FACULTAD DE ESTUDIOS SUPERIORES ACATLÁN
}

\vspace{0.15cm}

{\normalsize
Licenciatura en Matemáticas Aplicadas y Computación
}

% ---------------------------------------------------------
% ASIGNATURA
% ---------------------------------------------------------

\vspace{1.35cm}

{\large\bfseries
INVESTIGACIÓN EN MATEMÁTICAS APLICADAS Y\\
COMPUTACIÓN
}

% ---------------------------------------------------------
% TÍTULO
% ---------------------------------------------------------

\vspace{1.35cm}

{\LARGE\bfseries
Contenido sintético en Internet y degradación\\[0.15cm]
de modelos de lenguaje
}

\vspace{0.4cm}

{\large
Una aproximación al fenómeno de \textit{Model Collapse}
}

\end{center}

% ---------------------------------------------------------
% DATOS DEL TRABAJO
% ---------------------------------------------------------

% Espacio fijo en lugar de \vfill
\vspace{2.2cm}

\begin{flushright}

\begin{minipage}{0.58\textwidth}

\textbf{Alumnos:}\\[0.15cm]
Carlos Emiliano Segura Loera\\
Valeria Hernandez Santos\\
Nehiel Ramirez Hernandez 

\vspace{0.5cm}

\textbf{Profesor:}\\[0.1cm]
Sergio Quiroz Almaraz

\vspace{0.5cm}

\textbf{Fecha de entrega:}\\[0.1cm]
28 de septiembre de 2026

\vspace{0.5cm}

Naucalpan de Juárez, Estado de México

\end{minipage}

\end{flushright}

\end{spacing}

\end{titlepage}

% =========================================================
% PÁGINAS PRELIMINARES
% =========================================================

\pagenumbering{roman}
\setcounter{page}{1}

% =========================================================
% ÍNDICE
% =========================================================

\tableofcontents

\clearpage

% =========================================================
% RESUMEN
% =========================================================

\phantomsection
\addcontentsline{toc}{section}{Resumen}

\begin{center}

{\large\bfseries
RESUMEN
}

\end{center}

\vspace{0.5cm}

\justifying

El crecimiento de la inteligencia artificial generativa ha incrementado la
presencia de contenido sintético en Internet, incluyendo textos, imágenes,
publicaciones y otros materiales producidos total o parcialmente mediante
modelos generativos. Esta situación plantea interrogantes sobre las
consecuencias que puede tener la reutilización de dicha información en el
desarrollo de futuras generaciones de modelos de lenguaje. En particular, el
fenómeno conocido como \textit{model collapse} describe la degradación que
puede producirse cuando los modelos aprenden de manera recurrente a partir de
datos generados por otros modelos, favoreciendo patrones frecuentes y
reduciendo progresivamente la representación de información menos común. El
presente trabajo propone estudiar experimentalmente este fenómeno mediante
modelos de lenguaje pequeños capaces de ejecutarse localmente y conjuntos de
datos de alcance controlado. Se busca analizar los cambios producidos a través
de ciclos sucesivos de generación y reutilización de contenido sintético,
considerando principalmente la diversidad de las respuestas, la repetición de
patrones y las modificaciones en su distribución. Con ello se pretende
construir una aproximación accesible para observar y cuantificar algunos de
los efectos asociados con la degradación de modelos generativos.

\vspace{0.7cm}

\noindent
\textbf{Palabras clave:}
contenido sintético,
modelos de lenguaje,
\textit{model collapse},
datos sintéticos,
diversidad lingüística.

% CUERPO DEL ARTÍCULO EN DESARROLLO
% Los entregables de preguntas y justificación del capitulado se mantienen
% separados del artículo. Incorporar aquí únicamente contenido de investigación.
\clearpage
\pagenumbering{arabic}
\setcounter{page}{1}
\section{Introducción}

\subsection{Contexto y motivación}

La generación automática de texto plantea una cuestión relevante para el
desarrollo de modelos de lenguaje: qué ocurre cuando sus propias respuestas
se incorporan a los datos utilizados para entrenar modelos posteriores.
Si este proceso se repite, el contenido producido por una generación pasa
a formar parte de la información de la siguiente. En este escenario,
interesa estudiar si se conserva la variedad del contenido original o si
las respuestas se concentran progresivamente en determinados patrones.

El interés inicial de esta investigación surge de la llamada teoría del
``Internet muerto'', utilizada como punto de partida para reflexionar sobre
la presencia de contenido automatizado en la Web. Sin embargo, el estudio
no busca comprobar dicha teoría ni estimar qué proporción de Internet es
producida por inteligencia artificial. La motivación se concreta en un
problema experimental: analizar las consecuencias de reutilizar texto
sintético durante sucesivos ciclos de entrenamiento.

\subsection{Planteamiento del problema}

El fenómeno conocido como \textit{model collapse} describe un deterioro que
puede aparecer cuando los modelos se entrenan recursivamente con datos
generados por modelos anteriores. Shumailov et al.\ \cite{shumailov2024}
estudian este proceso y señalan la pérdida de información correspondiente
a regiones poco probables de la distribución original. Este antecedente
motiva observar tanto la variedad del texto generado como su concentración
en patrones frecuentes.

El efecto no debe suponerse idéntico en todas las condiciones. Kazdan
et al.\ \cite{kazdan2025} muestran diferencias entre reemplazar los datos
originales, acumularlos junto con datos sintéticos y entrenar con
subconjuntos de tamaño fijo. Por ello, la presencia de datos sintéticos no
es suficiente para anticipar el resultado: también importa el procedimiento
mediante el cual se incorporan al entrenamiento.

El problema de esta investigación consiste en determinar cómo cambia la
diversidad del contenido generado cuando se reutilizan datos sintéticos de
manera recurrente. En particular, se busca examinar los cambios léxicos y
semánticos, la repetición de patrones y su relación con la proporción de
datos sintéticos y el número de generaciones. El análisis deberá permitir
distinguir los cambios observados de una interpretación más amplia sobre
la aparición de efectos asociados con el \textit{model collapse}.

\subsection{Objetivos de la investigación}

El objetivo general es analizar experimentalmente los cambios en la
diversidad y repetición del contenido generado por un modelo de lenguaje
de escala reducida ante la reutilización recurrente de datos sintéticos.

De este objetivo se desprenden los siguientes objetivos específicos:
\begin{enumerate}
    \item Comparar la diversidad léxica y semántica del contenido generado
    a través de sucesivos ciclos de entrenamiento.
    \item Examinar cómo varían estas dimensiones al modificar la proporción
    de datos sintéticos utilizada.
    \item Evaluar los cambios en la repetición de patrones dentro de los
    textos y entre las respuestas generadas.
    \item Analizar el efecto conjunto de la proporción sintética y del
    número de generaciones sobre las dimensiones estudiadas.
    \item Examinar si los cambios observados son compatibles con efectos
    asociados con el \textit{model collapse}, considerando las condiciones
    y limitaciones del experimento.
\end{enumerate}

\subsection{Hipótesis}
El entrenamiento recurrente de modelos de lenguaje utilizando contenido generado por otros modelos produce una disminución progresiva en la diversidad del contenido generado y un incremento en la repetición de patrones lingüísticos. Se espera que estos efectos sean más pronunciados conforme aumenta la proporción de datos sintéticos empleada y el número de generaciones sucesivas de entrenamiento.

Esta hipótesis se contrastará dentro de las condiciones que se definan para
el experimento. No se presupone que el deterioro ocurra en todas las
proporciones ni que afecte del mismo modo a todas las dimensiones evaluadas.

\subsection{Justificación y aportación esperada}

Desde las matemáticas aplicadas y la computación, el problema permite
estudiar propiedades estadísticas del contenido generado y comparar su
evolución bajo condiciones controladas. Su interés radica en relacionar
la reutilización de datos con características observables del texto, sin
requerir reproducir el entrenamiento de un modelo comercial de gran escala.

La aportación esperada consiste en desarrollar una aproximación experimental
accesible y reproducible, con decisiones de entrenamiento y evaluación
documentadas. El fenómeno ya cuenta con antecedentes; por ello, la
contribución deberá precisarse a partir del modelo, los datos y el protocolo
seleccionados, así como del análisis de sus resultados.

\subsection{Alcances y delimitaciones}

El trabajo se centrará en generación de texto mediante modelos pequeños o
compactos y conjuntos de datos de tamaño controlado. No se contempla
entrenar modelos de gran tamaño desde cero. Se buscará utilizar datos
públicos o diseñados para el experimento, sin recurrir a información personal
ni a interacción con usuarios.

Se propone trabajar en español con un GPT-2 pequeño, datos de Wikipedia y
tres proporciones de objetivos sintéticos, durante cinco ciclos y tres
repeticiones. La metodología detalla esta configuración y las métricas
previstas. El piloto permitió comprobar el funcionamiento a escala reducida;
la ejecución del estudio principal sigue pendiente y dependerá de los recursos
computacionales y del tiempo disponible. Los hallazgos se interpretarán dentro de ese alcance,
sin extrapolarlos directamente a todos los modelos de lenguaje ni al
comportamiento de Internet en su conjunto.

\section{Marco teórico}

\subsection{Modelos generativos de lenguaje}

\subsubsection{Generación de texto como proceso probabilístico}

Un modelo de lenguaje asigna probabilidades a secuencias de unidades
textuales o \textit{tokens}. Estas unidades pueden corresponder a palabras,
fragmentos de palabras u otros símbolos, según la representación utilizada.
El modelado probabilístico busca aprender regularidades de las secuencias
observadas para estimar la probabilidad de otras secuencias
\cite{bengio2003}. En un modelo autorregresivo, la probabilidad de un texto
\(x=(x_1,\ldots,x_T)\) se expresa mediante la regla de la cadena:
\begin{equation}
    p_{\theta}(x)=\prod_{t=1}^{T}p_{\theta}(x_t\mid x_{<t}),
\end{equation}
donde \(\theta\) representa los parámetros aprendidos y \(x_{<t}\) el contexto
anterior al token \(t\). Generar texto consiste en seleccionar sucesivamente
tokens a partir de estas distribuciones condicionales.

La arquitectura Transformer introdujo un procesamiento de secuencias basado
en mecanismos de atención \cite{vaswani2017}. Este antecedente permite
comprender la arquitectura de numerosos modelos actuales, aunque no implica
que todos los modelos generativos tengan la misma configuración ni que en
este estudio ya se haya seleccionado uno.

\subsubsection{Preentrenamiento y ajuste de modelos}

El preentrenamiento permite aprender regularidades a partir de un corpus
inicial. El ajuste posterior parte de parámetros previamente aprendidos y
los modifica utilizando otros datos o un objetivo específico. Por ejemplo,
el ajuste con instrucciones utiliza ejemplos de solicitudes y respuestas
para modificar el comportamiento de un modelo preentrenado
\cite{wang2023selfinstruct}. Reentrenar desde cero y ajustar un modelo
existente son procedimientos distintos; esa diferencia deberá precisarse
al definir los ciclos experimentales.

\subsubsection{Estrategias de muestreo y decodificación}

La distribución aprendida no determina por sí sola el texto que se obtiene.
La decodificación voraz elige el token más probable, mientras que el muestreo
selecciona entre alternativas de manera probabilística. La temperatura
modifica la concentración de las probabilidades; \textit{top-k} restringe
la selección a un número de candidatos, y \textit{top-p} utiliza un conjunto
cuya masa de probabilidad acumulada alcanza un umbral. Holtzman et al.\
\cite{holtzman2020} muestran que la elección de decodificación puede cambiar
la diversidad y repetición sin modificar los parámetros del modelo.
Por ello, el análisis debe distinguir los efectos del entrenamiento de los
producidos por la generación.

\subsection{Datos utilizados en el entrenamiento}

\subsubsection{Datos de origen humano}

En este trabajo se denomina dato de origen humano al texto producido por
personas que se utilice como referencia inicial. Se trata de una definición
operativa: el origen no garantiza por sí mismo calidad, representatividad
ni ausencia de duplicados. La procedencia, las licencias y los criterios
de selección deberán documentarse para el corpus elegido.

\subsubsection{Datos sintéticos y semisintéticos}

Los datos sintéticos son ejemplos producidos mediante un modelo generativo.
Su utilización puede responder a finalidades diferentes. Self-Instruct
genera instrucciones y ejemplos, filtra casos inválidos o similares y los
utiliza para ajustar un modelo \cite{wang2023selfinstruct}. Los datos
semisintéticos combinan material original con transformaciones generadas;
Zhu et al.\ \cite{zhu2025} estudian la edición de tokens sobre textos humanos.
Por tanto, generar documentos completos y modificar parcialmente textos
existentes no constituyen tratamientos equivalentes.

\subsubsection{Mezcla, reemplazo y acumulación de datos}

Se distinguirán tres políticas: mezclar datos originales y sintéticos en
una proporción definida; reemplazar datos de una generación por los de la
siguiente; y acumular los conjuntos disponibles. Además, conservar un corpus
creciente y entrenar con todo él es distinto de extraer un subconjunto de
tamaño fijo en cada ciclo. Estas diferencias son relevantes para interpretar
el comportamiento del entrenamiento recursivo \cite{kazdan2025}. La
proporción sintética debe especificar su unidad de cálculo, por ejemplo
documentos o tokens, y no solo un porcentaje sin referencia.

\subsection{Entrenamiento recursivo}

\subsubsection{Ciclos de entrenamiento y generación}

El entrenamiento recursivo introduce una dependencia entre generaciones:
un modelo produce datos que se utilizan para entrenar o ajustar otro modelo.
Como esquema conceptual, puede representarse mediante
\begin{equation}
    D_g \longrightarrow M_g \longrightarrow S_g
    \longrightarrow D_{g+1},
\end{equation}
donde \(D_g\) es el conjunto usado en la generación \(g\), \(M_g\) el modelo
resultante y \(S_g\) sus datos sintéticos. El conjunto siguiente puede
incorporar datos originales y sintéticos, según la política seleccionada.
Este esquema describe la dependencia entre datos y modelos; no fija todavía
el protocolo experimental \cite{shumailov2024,kazdan2025}.

\subsubsection{Transmisión de información entre generaciones}

Cada conjunto generado es una muestra finita del comportamiento del modelo,
no una copia exacta de la distribución inicial. Un patrón poco probable
puede no aparecer en esa muestra y perder representación en los datos
posteriores. El ajuste puede añadir diferencias de aproximación respecto
de los datos recibidos. Shumailov et al.\ \cite{shumailov2024} analizan cómo
estos errores pueden acumularse. También deberá distinguirse si cada ciclo
parte de una copia del mismo modelo base o de los parámetros del ciclo
anterior.

\subsection{El fenómeno de \textit{model collapse}}

\subsubsection{Definición y manifestaciones}

El \textit{model collapse} se refiere a una degradación asociada con la
reutilización de datos generados por modelos anteriores. Puede manifestarse
como pérdida de información de la distribución original y reducción de
la variabilidad del contenido \cite{shumailov2024}. Su estudio requiere
describir el procedimiento y los indicadores observados, sin asumir que
todo uso de datos sintéticos tendrá el mismo efecto.

\subsubsection{Cambios en la distribución del contenido}

La velocidad del deterioro también requiere atención. Suresh et al.\
\cite{suresh2025} caracterizan tasas de colapso en distribuciones discretas
y gaussianas bajo estimación de máxima verosimilitud o aproximaciones a ella.
En esos escenarios simplificados, el olvido puede requerir muchas
iteraciones. Este resultado no fija una tasa para los modelos de lenguaje,
pero distingue la posibilidad de colapso de su rapidez observable en un
experimento de duración limitada.

La distribución de las respuestas puede cambiar en el vocabulario, los
patrones o los temas representados. Una mayor concentración implica que
ciertas alternativas reúnen una parte mayor de las observaciones. Las
frecuencias de palabras o de \(n\)-gramas permiten estudiar aspectos de ese
cambio, pero no describen por completo la distribución del lenguaje ni
su significado \cite{zhu2025}.

\subsubsection{Pérdida de patrones poco frecuentes}

Las colas de una distribución corresponden, en este contexto, a eventos
con baja probabilidad. Su pérdida puede preceder a una concentración más
severa del modelo \cite{shumailov2024}. Para estudiar patrones textuales
poco frecuentes será necesario definirlos respecto de un corpus de
referencia. Una reducción de vocabulario, por sí sola, no demuestra pérdida
de información semántica ni permite identificar qué casos se omitieron.

\subsubsection{Diferencias entre colapso, repetición y rendimiento}

La repetición es una propiedad observable del texto y puede aparecer por
la decodificación, sin entrenamiento recursivo \cite{holtzman2020}.
El rendimiento, por su parte, depende de una tarea o evaluación definida.
Drayson et al.\ \cite{drayson2025} consideran conjuntamente rendimiento,
características de generación y semejanza con referencias humanas. Estas
dimensiones pueden ofrecer información complementaria; ninguna se tomará
como prueba aislada y suficiente de colapso en este estudio.

\subsection{Diversidad del contenido generado}

\subsubsection{Diversidad léxica}

La diversidad léxica describe la variedad de vocabulario y combinaciones
de palabras. Dos conjuntos pueden diferir en términos distintos utilizados
o en la frecuencia de sus combinaciones. Su comparación exige considerar
el tamaño de muestra y las longitudes, pues afectan las medidas de variedad
\cite{guo2024}.

\subsubsection{Diversidad semántica}

La diversidad semántica se refiere a la variedad de significados expresados.
Cambiar palabras no implica necesariamente cambiar el contenido. Las
representaciones vectoriales de oraciones permiten aproximar relaciones
de similitud semántica; Sentence-BERT utiliza representaciones comparables
mediante similitud coseno \cite{reimers2019}. Esa aproximación depende del
modelo de representación, el idioma y la tarea, por lo que no debe tratarse
como una medición directa y completa del significado.

\subsubsection{Diversidad sintáctica}

La diversidad sintáctica corresponde a la variedad de estructuras con que
se organizan las oraciones. Guo et al.\ \cite{guo2024} distinguen las
dimensiones léxica, sintáctica y semántica, mostrando que no evolucionan
necesariamente con la misma intensidad. Se incluye esta dimensión como
fundamento conceptual; incorporarla a la evaluación dependerá de la
factibilidad y de las herramientas disponibles para el idioma elegido.

\subsection{Repetición y concentración de patrones}

\subsubsection{Repetición dentro de un texto}

Se considerará repetición interna la reiteración de palabras, secuencias
o fragmentos en una misma respuesta. Un \(n\)-grama es una secuencia de
\(n\) unidades consecutivas; sus frecuencias permiten describir
reiteraciones a distintas escalas. Se distinguirán las repeticiones
necesarias por el tema de aquellas que dominan la respuesta o forman
bucles, un comportamiento relacionado también con la decodificación
\cite{holtzman2020}.

\subsubsection{Similitud entre textos generados}

La similitud entre respuestas describe hasta qué punto distintos textos
repiten expresiones o contenidos. Un modelo puede generar respuestas
internamente variadas y, sin embargo, producir casi la misma respuesta
ante diferentes entradas. Texygen introduce Self-BLEU para comparar un
texto generado con otros del mismo conjunto \cite{zhu2018texygen}.
La similitud debe interpretarse junto con las entradas y la tarea: una
respuesta factual admite menos variación que una tarea creativa.

\subsection{Fundamentos de las métricas de evaluación}

\subsubsection{Medidas de diversidad léxica}

La razón de tipos a tokens, TTR, divide el número de palabras distintas
entre el total de palabras. Distinct-\(n\), utilizado en generación de
diálogo \cite{li2016}, relaciona las secuencias distintas de longitud
\(n\) con el total observado. Para un texto \(x\) con \(N\) unidades y
\(N\geq n\), se puede expresar como
\begin{equation}
    \operatorname{Distinct}\text{-}n(x)
    =\frac{|\operatorname{Unicos}_n(x)|}{N-n+1},
\end{equation}
donde el numerador cuenta los \(n\)-gramas distintos. Debe precisarse si
la medida se calcula por texto o sobre el conjunto concatenado. Guo
et al.\ \cite{guo2024} controlan la longitud al emplear estas medidas;
comparar tamaños diferentes puede confundir variedad y extensión.

\subsubsection{Medidas de diversidad semántica}

Para vectores de oración no nulos \(u\) y \(v\), la similitud coseno es
\begin{equation}
    \operatorname{sim}_{\cos}(u,v)
    =\frac{u\cdot v}{\lVert u\rVert\lVert v\rVert}.
\end{equation}
El promedio de \(1-\operatorname{sim}_{\cos}\) entre pares ofrece una
descripción de su separación en el espacio de representación. Su uso
para diversidad requiere fijar el modelo de representaciones y el
muestreo de pares, además de comprobar su adecuación al idioma
\cite{reimers2019,guo2024}.

\subsubsection{Medidas de repetición}

La fracción de \(n\)-gramas repetidos describe reiteración superficial
dentro de los textos. Self-BLEU toma cada texto como candidato y los demás
como referencias, de modo que un valor alto indica mayor semejanza entre
las salidas \cite{zhu2018texygen}. Es necesario especificar longitud del
\(n\)-grama, tokenización y tamaño del conjunto de referencia. Estas
medidas no identifican por sí solas equivalencia de significado.

\subsubsection{Medidas de fidelidad distribucional y rendimiento}

La perplejidad resume la probabilidad que un modelo asigna a una secuencia;
para el modelo autorregresivo anterior, se expresa como
\begin{equation}
    \operatorname{PPL}_{\theta}(x)
    =\exp\left(-\frac{1}{T}\sum_{t=1}^{T}
    \log p_{\theta}(x_t\mid x_{<t})\right).
\end{equation}
Su comparación requiere condiciones equivalentes de tokenización y datos
de evaluación. Una perplejidad baja no equivale necesariamente a mayor
diversidad o calidad del texto \cite{guo2024}.

MAUVE compara distribuciones de texto generado y humano mediante fronteras
de divergencia \cite{pillutla2021}. Ofrece información complementaria a
la variedad de palabras, pero no constituye una prueba automática de
corrección factual ni de colapso. La selección final de métricas se
realizará en la metodología; su presentación aquí establece los conceptos
para justificar esa selección y sus límites.

\section{Estado del arte}

\subsection{Criterios de búsqueda y selección de literatura}

La revisión se orientó a trabajos sobre entrenamiento recursivo con datos
sintéticos, diversidad lingüística y condiciones de mitigación. Se
consultaron fuentes primarias en Nature, ACL Anthology, PMLR y arXiv,
además de búsquedas por título y DOI en Scite. Se priorizaron antecedentes
que permiten comparar políticas de datos, generación de texto o evaluación.
Se distinguieron las publicaciones de sus preprints para evitar contar
versiones de un mismo trabajo como evidencia independiente.

Se trata de una revisión dirigida al diseño del estudio, no de una revisión
sistemática ni exhaustiva. El análisis combina resultados teóricos y
experimentos, precisando sus respectivos alcances. Los trabajos
seleccionados no utilizan necesariamente la misma definición operacional
de colapso ni los mismos indicadores; por ello, sus resultados requieren
compararse junto con el procedimiento que los produjo.

\subsection{Evidencia inicial y velocidad del deterioro}

Shumailov et al.\ \cite{shumailov2024} relacionan el entrenamiento recursivo
con la pérdida de información de la distribución original. Su evaluación
de lenguaje utiliza el ajuste de OPT-125m sobre WikiText-2 y generación
mediante búsqueda en haz. El trabajo constituye un antecedente para
observar cambios entre generaciones con modelos compactos, pero sus
condiciones no equivalen a cualquier uso posible de datos sintéticos.

Suresh et al.\ \cite{suresh2025} complementan esta perspectiva al estudiar
la velocidad de pérdida bajo familias de distribuciones simplificadas y
estimación de máxima verosimilitud o aproximaciones cercanas. Sus
resultados muestran que el tamaño de muestra influye en la escala temporal
del deterioro. La diferencia es relevante: demostrar una tendencia
asintótica no determina que un experimento corto con un modelo de lenguaje
muestre un colapso avanzado. Para el presente estudio, esto orienta a
describir la magnitud de los cambios sin exigir de antemano un desenlace
de colapso completo.

\subsection{Estudios sobre pérdida de diversidad lingüística}

Guo et al.\ \cite{guo2024} desplazan la atención desde el rendimiento hacia
la diversidad léxica, sintáctica y semántica en tareas de generación en
inglés. Mediante ajuste recursivo, encuentran reducciones especialmente
claras en las dimensiones léxica y sintáctica, mientras que la dimensión
semántica resulta relativamente más estable. Su evaluación incorpora
medidas de vocabulario, coincidencia de secuencias y representaciones
de oraciones, con controles de longitud para las métricas sensibles a ella.

Este antecedente se relaciona directamente con la pregunta de investigación:
el estudio de diversidad mediante modelos reducidos ya cuenta con evidencia
previa. En consecuencia, la aportación deberá precisarse en las condiciones
comparadas y en la documentación del procedimiento. También será necesario
justificar las herramientas lingüísticas si se utiliza un idioma diferente,
sin asumir que los resultados se trasladan automáticamente.

\subsection{Efectos de la proporción de datos sintéticos}

Zhu et al.\ \cite{zhu2025} estudian distintas proporciones sintéticas y
reportan una relación negativa con el rendimiento en sus configuraciones.
Su análisis identifica cambios distribucionales y concentración de
características de \(n\)-gramas. Además, proponen producir datos
semisintéticos mediante edición de tokens de textos humanos y evalúan
su uso en diferentes modalidades de entrenamiento.

Este trabajo aporta una razón para comparar proporciones, pero no establece
un porcentaje universal de deterioro. El origen, la construcción y la
selección del contenido intervienen junto con la cantidad utilizada.
Para el diseño propio, modificar una proporción exigirá explicitar qué
se conserva constante y qué se sustituye; de otro modo, el efecto podría
confundirse con diferencias de volumen o composición del corpus.

\subsection{Conservación de datos originales y acumulación}

Gerstgrasser et al.\ \cite{gerstgrasser2024} contrastan el reemplazo de datos
con la conservación de los originales y la acumulación de generaciones
sintéticas. En las condiciones evaluadas, la acumulación evita el deterioro
observado al reemplazar los datos. Su preprint introduce así una distinción
entre la presencia de contenido sintético y la política de conservación.

Kazdan et al.\ \cite{kazdan2025} comparan tres procedimientos: reemplazo,
acumulación con uso de todos los datos y acumulación con subconjuntos de
tamaño fijo. Reportan colapso con reemplazo, estabilidad de la pérdida de
prueba al usar toda la acumulación y deterioro gradual con subconjuntos
limitados. La lectura conjunta indica que conservar un corpus y utilizarlo
íntegramente son decisiones distintas. Además, estabilidad de pérdida de
prueba no demuestra por sí sola conservación de toda la diversidad.

\subsection{Influencia de las estrategias de generación}

Drayson et al.\ \cite{drayson2025} analizan la decodificación en experimentos
con GPT-2 y SmolLM2, considerando rendimiento, diversidad y semejanza con
referencias humanas. Sus comparaciones muestran que cambiar la estrategia
de generación modifica la trayectoria del deterioro. El procedimiento
que produce los datos sintéticos debe, por tanto, formar parte de la
descripción experimental, en lugar de tratarse como un detalle secundario.

Esta perspectiva también limita la interpretación de una puntuación de
diversidad aislada: una salida puede contener muchas secuencias diferentes
y aun así deteriorarse en otros aspectos. Para la investigación se propone
documentar entradas, longitudes y muestreo, y acompañar las medidas de
diversidad con una evaluación complementaria cuya viabilidad se determine
al cerrar el diseño.

\subsection{Estrategias de mitigación}

Los antecedentes revisados intervienen en distintos puntos del proceso.
Gerstgrasser et al.\ \cite{gerstgrasser2024} actúan sobre la conservación
del corpus. Zhu et al.\ \cite{zhu2025} modifican la construcción del dato
sintético al mantener parte del texto humano. Drayson et al.\
\cite{drayson2025} proponen remuestrear los datos utilizando estimaciones
de un detector para dar mayor presencia a contenido probablemente humano.

Estas intervenciones no constituyen una misma solución. Cada una introduce
supuestos y recursos propios: acceso al corpus original, posibilidad de
editarlo o capacidad de detectar el origen de las muestras. Su utilidad
para el proyecto será orientar controles y explicar diferencias entre
condiciones. La revisión no implica que todas deban implementarse ni que
sus mejoras estén garantizadas fuera de los escenarios estudiados.

\subsection{Comparación crítica de los enfoques}

La comparación permite identificar tres diferencias que deben conservarse
en el análisis. Primero, los resultados teóricos de distribuciones
simplificadas no fijan directamente una tasa para modelos de lenguaje.
Segundo, los estudios de acumulación pueden cambiar simultáneamente el
origen y el volumen de los datos. Tercero, las medidas de diversidad,
pérdida de prueba y calidad de generación evalúan dimensiones relacionadas,
pero distintas \cite{suresh2025,kazdan2025,guo2024,drayson2025}.

Por ello, diferencias entre hallazgos no deben interpretarse inmediatamente
como contradicciones. Antes se requiere comparar modelo, tarea, corpus,
inicialización de cada ciclo, política de conservación, presupuesto de
entrenamiento y decodificación. También es necesario precisar si las
repeticiones corresponden a cadenas de entrenamiento independientes o
solamente a varias respuestas de un mismo modelo.

\subsection{Delimitación de la aportación del estudio}

La literatura ya ha documentado tanto deterioro como condiciones que lo
modifican. La aportación prevista se sitúa en una reproducción o comparación
controlada y accesible que relacione proporción sintética y generaciones
con diversidad y repetición. Su valor dependerá de definir las condiciones,
documentarlas y analizar la incertidumbre de las observaciones, sin
presentar como nuevo el fenómeno ni prometer resultados todavía inexistentes.

Esta revisión orienta el diseño comparativo en español que se desarrolla
a continuación. La selección de recursos y controles se apoya también en
un piloto de viabilidad; ninguno de estos elementos permite anticipar el
sentido de los resultados de la comparación principal.

\section{Metodología}

\subsection{Enfoque y diseño experimental}

Se plantea un experimento comparativo en español con dos factores:
proporción sintética y número de generaciones. La configuración principal
se preparó después de un piloto de viabilidad, pero su ejecución completa
todavía está pendiente. Se mantendrán fijos el protocolo y los criterios de
evaluación antes de observar las comparaciones de prueba, sin exigir que
se produzca un deterioro o un colapso avanzado.

\subsection{Variables y condiciones de comparación}

La proporción sintética, \(\alpha\), se definirá respecto del número de
tokens objetivo utilizados para ajustar el modelo. Se proponen tres niveles:
\(\alpha=0\), \(0.5\) y \(1\). El nivel cero será la condición de referencia
con datos originales; el nivel intermedio conservará la mitad de los tokens
objetivo originales; el nivel uno los reemplazará por contenido sintético.
Se proponen cinco ciclos después de una generación inicial y tres cadenas
independientes por condición. Estos valores son objetivos de ejecución,
no cantidades verificadas experimentalmente.

Las variables de respuesta serán diversidad léxica, dispersión semántica,
repetición y perplejidad en datos reservados. Se mantendrán constantes
modelo base, tokenizador, número de ejemplos, longitudes, pasos de ajuste
y configuración de generación. La comparación no incluirá acumulación de
un corpus creciente; utilizará un presupuesto fijo por ciclo.

\subsection{Recursos e instrumentos computacionales}

El piloto se ejecutó en Google Colab con una Tesla T4. Para la ejecución
principal se prevé el mismo tipo de entorno, registrando Python, PyTorch,
Transformers, Datasets, Safetensors, CUDA y acelerador. Esa ejecución no
garantiza disponibilidad posterior de la GPU ni un tiempo fijo de servicio.
El notebook principal está preparado para guardar avances en una carpeta
persistente de Drive, cuyo montaje requiere autorización del usuario.
La reanudación comprobará que no cambien configuración, código o entorno
esencial dentro de una misma ejecución.

\subsection{Modelo y tarea propuestos}

Se utilizará \texttt{datificate/gpt2-small-spanish}, un GPT-2 pequeño
adaptado al español a partir de un checkpoint inglés mediante datos de
Wikipedia en español, según su ficha \cite{spanishmodelcard}. La tarea
será continuar fragmentos de texto en español, sin convertir el problema
en seguimiento de instrucciones. El modelo tiene 124\,439\,808 parámetros
en la carga verificada durante el piloto y permite un ajuste completo
de escala acotada. Su selección no implica reproducir exactamente los
experimentos de la literatura ni reclamar novedad por el idioma.
Se conservarán la revisión y las condiciones de uso del recurso.

\subsection{Selección y preparación de los datos}

Se utilizará \texttt{wikimedia/wikipedia}, configuración
\texttt{20231101.es}, con registros de artículos en español
\cite{wikipediacard}. La ficha ofrece una sola partición oficial;
entrenamiento, validación y prueba serán particiones propias por artículo.
Se aplicará SHA-256 a la semilla de partición y al ID del artículo, tomando
los primeros ocho bytes módulo 100: valores menores que 80 se asignarán
a entrenamiento, de 80 a 89 a validación y el resto a prueba. Se conserva
la semilla 20261109 del piloto para no trasladar sus artículos de ajuste
o validación a la referencia final de prueba.

Se recorrerá mediante streaming un prefijo de hasta 2000 artículos. El texto
se normalizará a Unicode NFC, conservando tildes y ñ, y se tokenizará sin
añadir tokens especiales. De cada artículo suficientemente largo se
seleccionará, con semilla fija, una ventana entre bloques no solapados de
64 tokens de contexto y 128 de objetivo. Después se seleccionarán 512
artículos para entrenamiento, 128 para validación y 128 para prueba.
Las tres muestras serán disjuntas y se comprobarán IDs, hashes de contenido
y ventanas exactas; una muestra insuficiente o duplicada detendrá la preparación.

Validación fijará la longitud común de las métricas antes de evaluar modelos
en prueba. Los 128 contextos de prueba y sus continuaciones originales serán
iguales para todas las condiciones y repeticiones. El prefijo acotado no es
representativo de toda Wikipedia ni de Internet. Además, el corpus no
certifica autoría humana individual y el modelo ya fue adaptado a Wikipedia:
la separación propia no demuestra ausencia de solapamiento con su entrenamiento
previo. Estas limitaciones se conservarán al interpretar la referencia y la
perplejidad.

\subsection{Presupuesto y definición de una generación}

La generación inicial, \(M_0\), se obtendrá ajustando el modelo base sobre
los 512 ejemplos originales. Para cada repetición, ese mismo punto inicial
servirá para las tres proporciones. Una generación posterior comprende:
producción de continuaciones con el modelo anterior, construcción de la
mezcla, ajuste y evaluación del modelo resultante.

Cada ajuste partirá de una copia del mismo checkpoint preentrenado, con
optimizador nuevo. Así, la transmisión entre generaciones ocurrirá mediante
los datos y no por la acumulación de actualizaciones en los parámetros.
Los contextos humanos permanecerán fijos en todos los niveles de mezcla;
la proporción sintética se aplicará exclusivamente a las continuaciones.
Por tanto, \(\alpha=1\) no significará ausencia total de contexto humano.

Cada conjunto contendrá \(512\times128=65\,536\) tokens objetivo. La mitad
sintética corresponderá a 256 continuaciones originales y 256 generadas,
seleccionadas mediante una permutación documentada. Los 256 IDs sintéticos
serán fijos a través de los ciclos de una misma repetición. Se generarán
512 continuaciones en cada ciclo no humano antes de aplicar la mezcla.
La comparación completa
requerirá 45 ajustes posteriores y tres ajustes iniciales, además del
muestreo y la evaluación. Este conteo deberá usarse al estimar el costo.

\subsection{Procedimiento de ajuste}

Se mantendrá una época por ajuste, tasa de aprendizaje \(2\times10^{-5}\),
AdamW con \(\beta_1=0.9\), \(\beta_2=0.999\), \(\epsilon=10^{-8}\)
y decaimiento de pesos 0.01 en todos los parámetros. Se utilizará float32,
recorte de norma de gradiente a 1, sin scheduler ni precisión mixta. El lote
efectivo será ocho mediante microbatches de dos y acumulación de gradientes.
La función de pérdida se calculará únicamente sobre los tokens objetivo;
los tokens del contexto y del relleno se enmascararán. Con 512 ejemplos y
lote efectivo ocho se obtendrán 64 actualizaciones por ajuste.

Se actualizarán todos los parámetros y se verificará la misma huella de
pesos preentrenados al iniciar cada ajuste, con optimizador vacío. Si se
requiere una modificación por recursos, se documentará en otra versión
del protocolo antes de continuar; no se mezclarán configuraciones ni se
seleccionarán checkpoints utilizando prueba.

\subsection{Generación de continuaciones}

Se utilizará temperatura \(1.0\) y \textit{top-p} \(0.95\), sin
\textit{top-k}, penalización de repetición ni búsqueda en haz. Se generarán
exactamente 128 tokens por continuación, excluyendo los IDs especiales,
incluido el de finalización. Esta restricción, verificada en el piloto,
modifica la distribución de decodificación y se mantendrá constante entre
condiciones. Los objetivos originales se tomarán sin agregar tokens
especiales. No se filtrarán salidas por idioma, calidad o valor de las
métricas; los textos y sus incidencias se conservarán.

La producción de datos del siguiente ciclo utilizará contextos de
entrenamiento; la evaluación utilizará los contextos reservados. Se
separarán las semillas de inicialización del ajuste, muestreo del corpus
y generación. Cada solicitud tendrá una semilla identificable, evitando
compartir y repetir inadvertidamente un mismo flujo aleatorio entre textos.
Las semillas del ajuste diferirán entre repeticiones y ciclos, y se
emparejarán entre proporciones del mismo ciclo. Las de evaluación serán
iguales por contexto dentro de una repetición a través de sus etapas.

\subsection{Métricas de evaluación}

Se calcularán Distinct-2 y Distinct-3 como n-gramas únicos entre n-gramas
totales del corpus, sin cruzar fronteras de texto. La repetición de
4-gramas será el promedio por salida de uno menos la proporción de
4-gramas únicos. Se utilizarán las primeras \(W\) palabras en minúsculas,
con \(W\) igual al mínimo entre 32 y la longitud de la continuación
original más corta de validación. Se conservarán tildes y ñ, sin stemming
ni eliminación de palabras vacías. Si alguna salida no alcanza \(W\), se
registrará y las métricas léxicas y semánticas de esa etapa quedarán
ausentes, sin excluir respuestas silenciosamente.

La diversidad semántica se aproximará mediante la media de la distancia
coseno entre pares de representaciones de las continuaciones, excluyendo
el contexto de entrada. Se utilizará
\texttt{sentence-transformers/paraphrase-multilingual-MiniLM-L12-v2},
codificador multilingüe que permite trabajar con español
\cite{multilingualcard}, con mean pooling enmascarado y normalización L2.
Su límite será 128 wordpieces, registrando truncamientos y usando la misma
revisión para todas las etapas. No se añadirá una métrica sintáctica en
este diseño de alcance reducido.

La perplejidad se calculará sobre los mismos tokens originales reservados,
con igual tokenización y enmascaramiento. Se acompañarán los indicadores
con los primeros ejemplos por orden guardado, para examinar coherencia
sin reemplazar las mediciones por una selección de casos llamativos.
La pérdida de patrones raros no se afirmará a partir de Distinct-n: su
evaluación específica requerirá un procedimiento adicional definido antes
del análisis y queda fuera del objetivo mínimo de ejecución.

\subsection{Análisis de resultados}

La unidad de réplica será la cadena de entrenamiento y no cada respuesta.
Se presentarán trayectorias por generación para cada proporción y
repetición. Se resumirán los cambios respecto de \(M_0\) y las diferencias
frente a la condición sin datos sintéticos en la misma generación.
Los contrastes de proporciones aprovecharán el punto inicial compartido
por repetición, manteniendo identificada esa dependencia.

Las métricas de prueba se guardarán por etapa, pero el notebook presentará
su resumen conjunto al completar el diseño; no se utilizarán para cambiar
hiperparámetros o decidir detención temprana. Con tres repeticiones se dará
prioridad a tamaños de cambio, medias y
rangos entre cadenas. Las comparaciones serán exploratorias; no se
presentarán las 128 respuestas como réplicas de entrenamiento independientes
ni se interpretará ausencia de significación como ausencia demostrada
de efecto. Si se modifica el número de repeticiones, se hará antes de
las comparaciones finales y se justificará la limitación correspondiente.

\subsection{Piloto y viabilidad temporal}

Se realizó un piloto separado con 64 ejemplos, 32 contextos de validación,
dos proporciones extremas, dos ciclos y una repetición. Completó cinco
ajustes y controles de enmascaramiento, generación y métricas; sus
resultados no se integrarán como réplicas de la comparación principal.
Los tiempos medidos orientan el presupuesto, pero no validan todavía la
duración de 48 ajustes ni su guardado persistente. Se prevé comprobar
primero una etapa principal y su reanudación, reservando tiempo para
análisis y redacción.

La estimación considerará todos los ciclos, condiciones y repeticiones,
además de guardado y posibles interrupciones. Si el diseño no cabe en el
plazo, se reducirá primero el número de ejemplos o generaciones, preservando
las tres condiciones y las repeticiones cuando sea posible. La modificación
se fijará por viabilidad, no por si el piloto muestra el efecto esperado.

\subsection{Reproducibilidad y consideraciones de uso}

Se conservarán configuración final, versiones, revisiones de modelos y
datos, manifiestos de particiones, semillas, registros de entrenamiento,
continuaciones y tablas de métricas. Las etapas completas se identificarán
mediante marcadores y hashes de archivos. Se guardarán lotes de generación
para continuar los pendientes; si se interrumpe un ajuste sin checkpoint
confirmado, se repetirá ese ajuste desde base y semilla. Se conservarán los
puntos iniciales y el checkpoint necesario para continuar, eliminando pesos
intermedios anteriores solo después de confirmar su sucesor. No se afirma
reanudación dentro de un optimizador ni conservación de todos los pesos.
No se utilizarán datos personales ni
se publicarán corpus o pesos sin revisar sus condiciones de uso. El informe
distinguirá lo propuesto de lo efectivamente ejecutado y explicará cualquier
desviación, incluida la imposibilidad de completar el diseño previsto.

% ESTRUCTURA DE TRABAJO: activar cada sección al incorporar su contenido.
% \section{Resultados}
% \section{Discusión}
% \section{Conclusiones}
% Trabajo futuro: opcional; puede resumirse al final de las conclusiones.

\clearpage
\phantomsection
\addcontentsline{toc}{section}{Referencias}
\begin{thebibliography}{99}

\bibitem{shumailov2024}
Shumailov, I., Shumaylov, Z., Zhao, Y., et al.
\textit{AI models collapse when trained on recursively generated data}.
Nature, 631, 755--759, 2024.
\url{https://doi.org/10.1038/s41586-024-07566-y}

\bibitem{kazdan2025}
Kazdan, J., Schaeffer, R., Dey, A., Gerstgrasser, M., Rafailov, R.,
Donoho, D. L., y Koyejo, S.
\textit{Collapse or Thrive: Perils and Promises of Synthetic Data in a
Self-Generating World}.
Proceedings of the 42nd International Conference on Machine Learning,
PMLR, 267, 29469--29494, 2025.
\url{https://proceedings.mlr.press/v267/kazdan25a.html}

\bibitem{bengio2003}
Bengio, Y., Ducharme, R., Vincent, P., y Jauvin, C.
\textit{A Neural Probabilistic Language Model}.
Journal of Machine Learning Research, 3, 1137--1155, 2003.
\url{https://jmlr.org/papers/v3/bengio03a.html}

\bibitem{vaswani2017}
Vaswani, A., Shazeer, N., Parmar, N., et al.
\textit{Attention Is All You Need}.
Advances in Neural Information Processing Systems, 30, 2017.
\url{https://papers.nips.cc/paper_files/paper/2017/hash/3f5ee243547dee91fbd053c1c4a845aa-Abstract.html}

\bibitem{holtzman2020}
Holtzman, A., Buys, J., Du, L., Forbes, M., y Choi, Y.
\textit{The Curious Case of Neural Text Degeneration}.
International Conference on Learning Representations, 2020.
Preprint inicial: arXiv:1904.09751, 2019.
\url{https://arxiv.org/abs/1904.09751}

\bibitem{wang2023selfinstruct}
Wang, Y., Kordi, Y., Mishra, S., Liu, A., Smith, N. A.,
Khashabi, D., y Hajishirzi, H.
\textit{Self-Instruct: Aligning Language Models with Self-Generated Instructions}.
Proceedings of ACL, 13484--13508, 2023.
\url{https://aclanthology.org/2023.acl-long.754/}

\bibitem{zhu2025}
Zhu, X., Cheng, D., Li, H., et al.
\textit{How to Synthesize Text Data without Model Collapse?}.
Proceedings of ICML, PMLR, 267, 79746--79771, 2025.
\url{https://proceedings.mlr.press/v267/zhu25d.html}

\bibitem{drayson2025}
Drayson, G., Yilmaz, E., y Lampos, V.
\textit{Machine-generated text detection prevents language model collapse}.
Proceedings of EMNLP, 29657--29673, 2025.
\url{https://aclanthology.org/2025.emnlp-main.1506/}

\bibitem{guo2024}
Guo, Y., Shang, G., Vazirgiannis, M., y Clavel, C.
\textit{The Curious Decline of Linguistic Diversity:
Training Language Models on Synthetic Text}.
Findings of NAACL, 3589--3604, 2024.
\url{https://aclanthology.org/2024.findings-naacl.228/}

\bibitem{reimers2019}
Reimers, N., y Gurevych, I.
\textit{Sentence-BERT: Sentence Embeddings using Siamese BERT-Networks}.
Proceedings of EMNLP-IJCNLP, 3982--3992, 2019.
\url{https://aclanthology.org/D19-1410/}

\bibitem{li2016}
Li, J., Galley, M., Brockett, C., Gao, J., y Dolan, B.
\textit{A Diversity-Promoting Objective Function for Neural Conversation Models}.
Proceedings of NAACL-HLT, 110--119, 2016.
\url{https://aclanthology.org/N16-1014/}

\bibitem{zhu2018texygen}
Zhu, Y., Lu, S., Zheng, L., Guo, J., Zhang, W., Wang, J., y Yu, Y.
\textit{Texygen: A Benchmarking Platform for Text Generation Models}.
Proceedings of SIGIR, 1097--1100, 2018.
\url{https://doi.org/10.1145/3209978.3210080}

\bibitem{pillutla2021}
Pillutla, K., Swayamdipta, S., Zellers, R., Thickstun, J.,
Welleck, S., Choi, Y., y Harchaoui, Z.
\textit{MAUVE: Measuring the Gap Between Neural Text and Human Text
using Divergence Frontiers}.
Advances in Neural Information Processing Systems, 34, 2021.
\url{https://papers.nips.cc/paper_files/paper/2021/hash/260c2432a0eecc28ce03c10dadc078a4-Abstract.html}

\bibitem{suresh2025}
Suresh, A. T., Thangaraj, A., y Khandavally, A. N. K.
\textit{Rate of Model Collapse in Recursive Training}.
Proceedings of the 28th International Conference on Artificial Intelligence
and Statistics, PMLR, 258, 1396--1404, 2025.
\url{https://proceedings.mlr.press/v258/suresh25a.html}

\bibitem{gerstgrasser2024}
Gerstgrasser, M., Schaeffer, R., Dey, A., et al.
\textit{Is Model Collapse Inevitable? Breaking the Curse of Recursion
by Accumulating Real and Synthetic Data}.
Preprint arXiv:2404.01413v2, 2024.
\url{https://doi.org/10.48550/arXiv.2404.01413}

\bibitem{spanishmodelcard}
Datificate. \textit{Model card: gpt2-small-spanish}.
Ficha del modelo consultada el 9 de noviembre de 2026.
\url{https://huggingface.co/datificate/gpt2-small-spanish}

\bibitem{wikipediacard}
Wikimedia. \textit{Dataset card: Wikipedia}, configuración 20231101.es.
Ficha del conjunto consultada el 9 de noviembre de 2026.
\url{https://huggingface.co/datasets/wikimedia/wikipedia}

\bibitem{multilingualcard}
Sentence Transformers. \textit{Model card: paraphrase-multilingual-MiniLM-L12-v2}.
Ficha del codificador consultada el 9 de noviembre de 2026.
\url{https://huggingface.co/sentence-transformers/paraphrase-multilingual-MiniLM-L12-v2}

\end{thebibliography}

\end{document}
